\appendix
\section*{Appendix A. DCE Instrument, Choice Tasks, and Validity Diagnostics}
\addcontentsline{toc}{section}{Appendix A. Internal Validity Diagnostics and DCE Materials}

This appendix presents the diagnostic checks used to assess internal validity in the discrete 
choice experiment (DCE), including (i) dominance behaviour, (ii) attribute non-attendance 
(ANA), (iii) transitivity and consistency, (iv) intra-individual variable correlations, and 
(v) pilot-study procedures. Collectively, these checks support the internal reliability and 
behavioural coherence of the responses.

\subsection*{A.1 Dominance Behaviour}

Dominance tests evaluate whether respondents ever selected an alternative that was strictly 
dominated across all attributes. As shown in Table~A1, only 1.8\% of respondents violated 
dominance at least once—well within the $\leq 5\%$ threshold commonly considered acceptable 
in DCE research \cite{Louviere2000,Johnson2013}. This low rate suggests that respondents 
understood the choice tasks and engaged in meaningful trade-offs.

\subsection*{A.2 Attribute Non-Attendance (ANA)}

ANA analyses assess whether respondents systematically ignored particular attributes during 
choice tasks. Table~A1 indicates that ANA was 0\% for vaccine delay, efficacy, side-effects, 
and cash incentives, reflecting strong salience of the primary decision attributes. A small 
minority (5.4\%) showed insensitivity to vaccine origin, suggesting minor simplification but 
no evidence of widespread attribute non-attendance. Overall, ANA diagnostics demonstrate 
robust engagement with the full attribute structure.

\begin{table}[H]
    \centering
    \caption{Dominance behaviour and attribute non-attendance (ANA)}
    \label{tab:validity}
    \begin{tabular}{lcc}
        \toprule
        \textbf{Check} & \textbf{Definition} & \textbf{Result (\%)} \\
        \midrule
        Dominance rate &
        Share selecting a dominated alternative &
        1.8 \\
        ANA: Vaccine delays &
        Choices invariant to delay levels &
        0.0 \\
        ANA: Vaccine efficacy &
        Choices invariant to efficacy levels &
        0.0 \\
        ANA: Side effects &
        Choices invariant to side-effect levels &
        0.0 \\
        ANA: Cash incentives &
        Choices invariant to cash levels &
        0.0 \\
        ANA: Vaccine origin &
        Choices invariant to origin levels &
        5.4 \\
        \bottomrule
    \end{tabular}

    \vspace{4pt}
    \footnotesize
    \textit{Note.} Dominance reflects selecting an option strictly inferior across all attributes. ANA 
    indicates choices unaffected by variation in a given attribute across tasks.
\end{table}

\subsection*{A.3 Transitivity and Choice Consistency}

Respondents’ choices satisfied the transitivity assumptions of random utility theory. No 
systematic violations or preference reversals were detected, and repeated patterns across the 
six tasks reflected stable and ordered preferences. These findings support the internal 
consistency and construct validity of the DCE, reinforcing the reliability of the resulting 
preference and time-discount parameter estimates.

\subsection*{A.4 Correlation Matrix of Intra-Individual Variables}

Figure~A1 displays pairwise correlations across all individual-level covariates included in 
the mixed-logit and subgroup models. All coefficients were below 0.10, indicating minimal 
concerns regarding multicollinearity or redundancy.

Access to healthcare showed only weak associations with other variables. Government trust was 
slightly positively correlated with age and rural residence, while smoking exhibited a weak 
negative correlation with chronic medical conditions. Overall, these low correlations confirm 
sufficient distinctness across covariates for use in the multivariate analyses.

\begin{figure}[H]
    \centering
    \includegraphics[width=0.9\textwidth]{correlation among Individual variables4.png}
    \caption{Correlation matrix of individual-level covariates. Darker colours indicate stronger 
    correlations. All coefficients are below 0.10, confirming negligible multicollinearity.}
    \label{fig:correlation}
\end{figure}

\subsection*{A.5 Example Choice Task}

Figure~A2 provides an example of the choice task presented to respondents. Each task 
featured two hypothetical vaccine profiles varying across five attributes—delay to 
vaccination, origin, efficacy, side-effect severity (supported by illustrative icons), and 
cash incentives—and an opt-out alternative. Attribute descriptions appeared at the top of the 
screen. For example, efficacy was defined behaviourally as the percentage reduction in 
infection risk (``95\% effectiveness means a 95\% reduction in infection risk''), and side 
effects were described using both text and images to reduce cognitive burden.

The example shown contrasts two vaccines identical in origin and efficacy but differing in 
delay (0 vs.\ 1 month), side-effect severity (fatigue for one week vs.\ fever lasting two 
days), and incentives (800 CNY vs.\ 200 CNY).

\begin{figure}[H]
\centering
\includegraphics[width=0.95\textwidth]{sample card.png}
\caption{Example DCE choice task shown to respondents. Participants selected between two hypothetical 
vaccine profiles or an opt-out option.}
\label{fig:sampletask}
\end{figure}

\subsection*{A.6 Pilot Study Procedures}

A pilot study was conducted with 60 adults from Wuhan to assess the clarity, feasibility, and 
cognitive burden of the DCE instrument. The pilot followed the same screening criteria as the 
main study and included the full sequence of choice tasks, sociodemographic questions, and a 
debriefing module.

The pilot served three purposes. First, it evaluated comprehension of attribute descriptions 
and level wording. Respondents demonstrated clear understanding of delay, efficacy, and 
side-effect attributes, while minor revisions were made to simplify phrasing for vaccine 
origin and severity icons.

Second, the pilot assessed task difficulty and respondent burden. Completion times ranged 
from $4.7$ to $10.2$ minutes (mean: $6.1$)
, indicating that the number of tasks and cognitive 
demands were appropriate. Dominance and ANA rates were consistent with those in the main 
sample, suggesting no evidence of satisficing or disengagement.

Third, the pilot refined the statistical efficiency of the experimental design. Preliminary 
parameter estimates informed updated priors in the efficient design algorithm, improving 
attribute balance and orthogonality in the final DCE.

Overall, pilot results confirmed that the survey instrument was clear, feasible, and capable 
of generating high-quality preference data.
